\chapter[All About Posture]{Proper Posture for a Happy and Healthy Guitar Life}

In this chapter, we will discuss the fundamental importance of playing the guitar with proper posture. Posture may not be the most glamorous or exciting aspect of playing guitar, but it is the thing you must understand first to be able to do all of the glamorous and exciting things later.

\section{Seated posture - getting the body set}

Everything starts with a chair. Make sure the chair you use has the following characteristics:

\begin{itemize}
	
	\item Made of a sturdy material like wood or plastic. Avoid chairs with an excessive amount of cushion.
	
	\item Allows you to sit with your hips and knees more or less on the same plane.
	
	\item Lets your feet touch the floor when you sit at the front edge. Do not sit far back in the chair seat. 
	
\end{itemize}  

Once you have an appropriate chair, practice the following sequence of actions. Be mindful of each step, and repeat this process multiple times.

\begin{enumerate}
	
	\item Sit at the front edge of the chair. Imagine feeling a cool breeze behind your back. Do not let your spine touch the back of the chair.
	
	\item Check that both feet are flat on the floor, or at least your toes and the front half of your feet are touching the floor. It's okay if your heels are slightly off of the floor.
	
	\item Place your hands on your thighs with your fingers pointed towards each other. Your elbows should be pointing out, and you should feel a large space inside of each arm.
	
	\item Try to raise your chest and head a little higher. Imagine that there is a string tied to the top of your head and you are being pulled gently upwards.
	
	\item Hold this posture for at least 10 seconds before releasing. Repeat multiple times. 
	
\end{enumerate}

\section{Adding the footstool}

Something that distinguishes the guitar from other instruments is the use of a footstool to raise the left foot off the floor. This, in turn, elevates the left leg and raises the neck of the guitar higher so that the strings and frets are closer to the view of the player. This also makes shifting to higher frets easier because gravity helps you do the work.

% Image of footstool


Look down at the floor at your left foot. Remember that spot on the floor, then move your foot and replace it with the footstool set to the lowest height setting. Footstools are typically height-adjustable, but you should always start with the lowest setting and only move higher if the guitar does not feel sufficiently elevated. Try to rest your foot lightly on the footstool. If you put too much of your body weight on the footstool, it may fall over.

\section{The four main contact points}

The use of the footstool also allows for a very secure posture where the guitar touches the body in four primary places.

\begin{enumerate}
	
	\item The left leg fits into the lower curve of the guitar's body. Make sure the neck of the guitar is pointing to the left.
	
	\item The right thigh touches the bottom right ``corner" of the guitar. Imagine the guitar body as rectangular and place the bottom right corner on your right thigh.
	
	\item The chest touches the top left ``corner" of the guitar. Again, imagine the guitar body as rectangular and left the top left corner touch your chest.
	
	\item The right forearm touches the top right ``corner" of the guitar. Let the forearm touch above the saddle and let ody of the guitarthe right hand hang down over the sound hole.
	
\end{enumerate}

% Image of contact points on guitar

\section{A word about tension and relaxation}

\section{Finding balance in every action}



